
Evaluating the factuality of LLMs is pivotal for ensuring the reliability and trustworthiness of their generated content \citep{pezeshkpour2023measuring,lee2022factuality}. As LLMs become increasingly integrated into various applications, from information retrieval to content generation, the accuracy of their outputs becomes paramount. In this section, we delve into the evaluation metrics and benchmarks used for assessing the factuality of LLMs, studies that have undertaken such evaluations, and domain-specific evaluation.
% \footnote{We've observed several studies focusing on Factuality Evaluation and Domain-Specific LLMs, including XX \citep{} and YY \citep{}. However, since they haven't published their findings, we've excluded them from this survey.}



\subsection{Factuality Evaluation Metrics}
\label{sec:eval_metrics}
\begin{table*}[t]\small
\centering
\caption{Evaluation Metrics for the Factuality of LLMs.}
\vspace{-2mm}
\label{tab:eval-metrics}
\begin{adjustbox}{max width=\textwidth}
  \setlength{\tabcolsep}{.4mm}
{
\begin{tabular}{l|l|l}
\toprule
\textbf{Evaluation Type} & \textbf{Metric} & \textbf{Description, Purpose/Usage, and Reference/Findings} \\ 
\midrule
\multirow{10}{*}{Rule-based} & Exact Match & Correct matches between generated and input/reference text. \\ \\
& \begin{tabular}[c]{@{}l@{}}Common Metrics \\ (Accuracy, Precision,\\
Recall,  F-Measure,\\ Calibration, Brier)\end{tabular} & \begin{tabular}[c]{@{}l@{}}\textbf{Accuracy}: Measures overall correctness of model prediction.\\ \textbf{Precision}: Ability of model to classify instances positively.\\ \textbf{Recall}: Ability of model to identify all positive instances.\\ \textbf{F-Measure}: Combines both precision and recall into a single score.\\ \textbf{Calibration}: Agreement between predicted probabilities and observed frequencies.\\ \textbf{Brier}: Measures accuracy of probabilistic predictions.\end{tabular} \\ \\
& MC1 and MC2 \citep{TruthfulQA} & Used to identify the correct answer from multi-choice questions. \\ 
& BLEU \citep{papineni2002bleu} & Determines frequency of phrase co-occurrence in two sentences. \\ 
& ROUGE \citep{lin2004rouge}& Calculates similarity between reference and generated text. \\ 
& METEOR\citep{lowe2017towards} & Comprehensive measure based on harmonic mean of unigram precision and recall. \\  \\
& QUIP-Score \citep{weller2023according}& Calculates how much of generated text consists of exact spans found in a text corpus. \\ 
\midrule
\multirow{5}{*}{Neural} & ADEM \citep{lowe2017towards} & \begin{tabular}[c]{@{}l@{}}Uses a Hierarchical RNN to evaluate the quality of responses in language models.\end{tabular}  \\
& BERTScore \citep{zhang2020bertscore}& \begin{tabular}[c]{@{}l@{}}Pre-trained embeddings from BERT used to evaluate sentence similarity.\end{tabular} \\ 
& BLEURT \citep{sellam2020bleurt}& \begin{tabular}[c]{@{}l@{}}Trains BERT on larger dataset and rate similarity of sentences .\end{tabular} \\ 
& BARTScore  \citep{yuan2021bartscore}& \begin{tabular}[c]{@{}l@{}}Evaluates quality using pre-trained sequence-to-sequence models.\end{tabular} \\ 
\midrule
%\hline
\multirow{2}{*}{Human} & \begin{tabular}[c]{@{}l@{}}AIS \citep{rashkin2023measuring}, \\Auto-AIS\citep{ gao2023rarr}\end{tabular} & \begin{tabular}[c]{@{}l@{}}Evaluates if output is backed by evidence.\end{tabular} \\ \\
& FActScore \citep{Min2023-FActScore}& \begin{tabular}[c]{@{}l@{}}Measures factual accuracy of LLMs breakpointing generated content into atomic facts.\end{tabular} \\ \midrule
\multirow{7}{*}{LLM-based} & GPTScore \citep{fu2023gptscore} & \begin{tabular}[c]{@{}l@{}}Evaluates quality of AI output. It is efficient and avoids the annotation requirement.\end{tabular} \\ \\
& GPT-judge  \citep{TruthfulQA}& \begin{tabular}[c]{@{}l@{}}Evaluates truthfulness of LLM answers.  It is used in evaluating truthfulness\\ and informativeness.\end{tabular} \\ \\
& \begin{tabular}[c]{@{}l@{}}Truthfulness and\\ Informativeness \\\citep{TruthfulQA}\end{tabular} & \begin{tabular}[c]{@{}l@{}}Truthfulness Measures the honesty of LLM information. \\ It is used to evaluate the factuality of information.\\
Informativeness Evaluates the relevance and value of  LLM responses 
\end{tabular}\\ \\
& LLM-Eval \citep{lin2023llm} & \begin{tabular}[c]{@{}l@{}}Evaluates the quality of a conversation. It is adaptable in various scenarios.\end{tabular} \\ \bottomrule
%\hline
\end{tabular}
}
\end{adjustbox}
\end{table*}


In this subsection, we delve into metrics established for evaluating the factuality of LLMs. As the problem formulation is akin to natural language generation (NLG)
~\citep{celikyilmaz2021evaluation,ji2023survey}, we introduce several automatic evaluation metrics typically used for NLG, as well as specifically examining the metrics for factuality.

We categorize these metrics into the following groups:
(1) Rule-based evaluation metrics;
(2) Neural evaluation metrics;
(3) Human evaluation metrics; and
(4) LLM-based evaluation metrics.

We list those metrics in Table \ref{tab:eval-metrics}.
\subsubsection{Rule-based evaluation metrics}

Most assessments of factuality in large language models use rule-based evaluation metrics, due to their consistency, predictability, and ease of implementation; they allow for reproducible outcomes through a systematic method. However, they can be rigid and may not account for nuances or variations in language use, context interpretation, or colloquial expressions. This means language models rated highly by these metrics may still produce content that feels unnatural or inauthentic to human readers.


\paragraph{Exact Match} An ``exact match" refers to a situation where the generated text precisely matches a specific input or reference text. This means that the LLM produces output that is identical, word-for-word, to the provided input or reference text. Exact matches are often used in NLG when you want to replicate or repeat a piece of text without any variations or alterations. Exact Match is commonly used in Open-domain Question Answering \citep{Fusion-in-Decoder,rfid}.

\paragraph{Common metrics}
Many factuality evaluation measurements use commonly adapted metrics, such as Accuracy, Precision, Recall, AUC, F-Measure, Calibration score,  Brier score, and other common metrics used in probabilistic forecasting and machine learning, specifically in tasks involving probabilistic predictions. The common definition of these metrics involves the use of correctly predicted labels and ground-truth labels. As the input and output of LLMs are human-readable sentences, there is no unified method to convert the sentences into labels. Most evaluation will define their own way. That is, these scores are frequently not used in isolation, but rather in combination. For instance, BERTScore~\citep{BERTscore} uses BERT for determining the Precision and Recall, then uses F-measures to get a final weighted score. In the following, we will describe the most simple form of these scores.


The \emph{Calibration score}, used in  \citep{lin2022teaching, kadavath2022language} measures the agreement between predicted probabilities and observed frequencies. A perfectly calibrated model should, over a large number of instances, see the predicted probability of an outcome match the relative frequency of that outcome.

The \emph{Brier Score}, used in ~\citep{kadavath2022language} is a metric used in probabilistic forecasting to measure the accuracy of probabilistic predictions. It calculates the mean squared difference between the predicted probability assigned to an event and the actual outcome of the event. 
The Brier Score ranges from 0 to 1, where 0 indicates a perfect prediction and 1 indicates the worst possible prediction. In other words, the lower the Brier Score, the better the accuracy of the prediction. It's worth noting that this metric is appropriate for binary and categorical outcomes, but not for ordinal outcomes. 
For binary outcomes, the Brier Score can be calculated as follows:

\begin{equation}
BS = \frac{1}{N} \sum_{i=1}^{N} (forecast_{i} - actual_{i})^2
\end{equation}

where \(forecast_{i}\) is the predicted probability, \(actual_{i}\) is the actual outcome (0 or 1), and \(N\) is the total number of forecasts made.


\paragraph{MC1 (Single-true) and MC2 (Multi-true)} are widely recognized metrics in multi-choice question answering, particularly in TruthfulQA \citep{TruthfulQA}.
MC1: For a given question accompanied by several answer choices, the objective is to identify the sole correct answer. The model's selection is determined by the answer choice to which it allocates the highest log probability of completion, independent of the other choices. The score is calculated as the straightforward accuracy across all questions.
MC2: Presented with a question and multiple reference answers labeled as true or false, the score is derived from the normalized total probability assigned to the collection of true answers.

\paragraph{BLEU}~\citep{papineni2002bleu}, also known as Bilingual Evaluation Understudy metric is commonly employed in the context of factuality evaluation. This metric calculates the co-occurrence frequency of phrases in two sentences, based on a weighted average of matched n-gram phrases. This helps in quantitatively assessing the factual consistency between the generated text and its reference.
% The BLEU score is computed as 

% \begin{equation}
% BLEU = BP * \exp\left(\frac{1}{N} * \sum_{n=1}^{N} P_n\right)
% \end{equation}

% \noindent where 
% \emph{(1)} \(BP\) (Brevity Penalty) accounts for the length of the candidate translation. If the candidate translation is shorter than the reference, \(BP\) is less than 1, which reduces the BLEU score.
% \emph{(2)} \(P_n\) is the n-gram precision, which measures the portion of n-grams in the candidate translation that are also present in the reference translation.
% \emph{(3)} \(N\) is the maximum order of n-grams considered (commonly up to 4 in much of the literature) and \(\sum_{n=1}^{N} P_n\) is the sum of \(\log P_n\) for \(n\) from 1 to \(N\).
% - \(\exp\) denotes the exponential function.

\paragraph{ROUGE}
The Recall-Oriented Understudy for Gisting Evaluation (ROUGE) metric \citep{lin2004rouge} serves as a measure of the similarity between the generated text and the reference text, with the similarity grounded on recall scores. Primarily used in the field of text summarization, the ROUGE metric incorporates four distinct types. These include ROUGE-n, which assesses n-gram co-occurrence statistics, and ROUGE-l, which measures the longest common subsequence. ROUGE-w provides evaluation based on the weighted longest common subsequence, while ROUGE-s measures skip-bigram co-occurrence statistics. These diverse metrics collectively provide a comprehensive measure of the factual accuracy of generated text.
% ROUGE score can be calculated in various ways based on the length of n-grams (unigram, bigram, etc.) used. The simplest version, ROUGE-n score can be compactly represented as:

% \begin{equation}
% \text{ROUGE-n} = \frac{\sum_{S \in RS}\sum_{gram_n \in S} Count_{match}(gram_n)}{\sum_{S \in RS}\sum_{gram_n \in S} Count(gram_n)}
% \end{equation}

% \noindent where:
% \emph{(1)} \(RS\) denotes the set of Reference Summaries,
% \emph{(2)} \(gram_n\) represents an n-gram within the reference summary \(S\),
% \emph{(3)} \(Count_{match}(gram_n)\) signifies the number of times that n-gram \(gram_n\) appears in both the generated text and the reference summary,
% \emph{(4)} \(Count(gram_n)\) is the frequency of occurrence of the n-gram \(gram_n\) within the reference summary.


\paragraph{METEOR}
The Metric for Evaluation of Translation with Explicit Ordering (METEOR) \citep{banerjee2005meteor} aims to address several shortcomings presented by BLEU. These include deficiencies in recall, the absence of higher order n-grams, an absence of explicit word-matching between the generated and reference text, and the use of geometric averaging of n-grams. METEOR overcomes these by introducing a comprehensive measure, calculated based on the harmonic mean of the unigram precision and recall. This offers potentially enhanced appraisals of factuality in the generated text.


\paragraph{QUIP-Score}~\citep{weller2023according}  is an n-gram overlap measure. It quantifies the degree to which a generated passage consists of exact spans found in a text corpus. The QUIP-Score serves to evaluate the 'grounding' ability of LLMs, specifically assessing whether model-generated answers can be directly located within the underlying text corpus. It is defined by comparing the precision of the character n-gram from the generated output to the pre-training corpus. 
% This is formally illustrated by generation $Y$ and the text corpus $C$:
% \begin{equation}
% \operatorname{QUIP}(Y ; C)=\frac{\sum_{\operatorname{gram}_n \in Y} \mathbb{1}_C\left(\operatorname{gram}_n\right)}{\left|\operatorname{gram}_n \in Y\right|},
% \end{equation}
% where $\mathbb{1}($.$) $is an indicator function:
% \begin{equation}
% \mathbb{1}(.) = \begin{cases}
%     1, & \text{if } \text{gram} \in C\\
%     0, & \text{otherwise}
%     \end{cases}
% \end{equation}

% Therefore, a score of 0.5 implies that 50\% of the n-grams derived from the generated text can be found in the pre-training corpus. The authors calculate a macro-average of this value over a collection of generations, which results in a single performance figure representative of a specific test dataset.




\subsubsection{Neural  Evaluation Metrics}
These metrics operate by comparing the output of these models with a standard or reference text by learning evaluator models. This category primarily comprises three prominent metrics, namely ADEM, BLEURT, and BERTScore. Each metric approaches the evaluation slightly differently, yet all aim to assess the semantic and lexical alignment between machine-generated text and its reference counterpart, thus ensuring the factuality of the generated content.

\paragraph{ADEM}
The automatic Dialogue Evaluation Model (ADEM) metric is a cogent tool utilized to gauge the quality of responses generated by language models in a conversation. Developed by \citet{lowe2017towards}, this metric trains a Hierarchical RNN, in a semi-supervised fashion, to predict ratings for the machine-generated responses. The ADEM's assessment primarily hinges on a dialogue context, designated as $\mathbf{c}$, alongside the model's response, labeled as $\mathbf{\hat{r}}$, and a reference response, specified as $\mathbf{r}$. These elements, when encoded via the hierarchical RNN, inform the ADEM which then predicts a score, reflecting the proximity of the model's response to factual accuracy and relevance.
% :
% \begin{equation}
% score = (\mathbf{c}^{\top} M \mathbf{\hat{r}} + \mathbf{r}^{\top} N \mathbf{\hat{r}} - \alpha) / \beta,
% \end{equation}
% where $M, N$ are learnable parameters and $\alpha, \beta$, represent scalar constants that serve as initial values. 




\paragraph{BERTScore}
the BERTScore metric, as introduced by \citet{zhang2020bertscore}, utilizes pre-trained embeddings from BERT to gauge the similarity between two sentences. This is done by assigning embeddings to a referenced sentence, "x", and a model-generated sentence, "$\hat{x}$", denoted as "$\mathbf{x}$" and "$ \mathbf{\hat{x}}$", respectively. The recall, precision, and F1-scores are then calculated to quantify the similarity between "$\mathbf{x}$" and "$ \mathbf{\hat{x}}$". This score provides a measure of the generated sentence's factuality.
% , and thus, the reliability of the large language model itself:
% \begin{align}
% R_{\rm BERT} &= \frac{1}{|x|} \sum_{x_i \in x} \max_{\hat{x}_j \in \hat{x}} \mathbf{x}_i^{\top} \mathbf{\hat{x}}_j, \\
% P_{\rm BERT} &= \frac{1}{|\hat{x}|} \sum_{\hat{x}_j \in \hat{x}} \max_{x_i \in x} \mathbf{x}_i^{\top} \mathbf{\hat{x}}_j, \\
% F_{\rm BERT} &= 2 \frac{P_{\rm BERT} \cdot R_{\rm BERT}}{P_{\rm BERT} + R_{\rm BERT}},
% \end{align}
% where the recall and precision elements are calculated based on a token-match approach. The recall score stems from comparing each token in the reference sentence "$x$" to the corresponding token in the generated sentence "$\hat{x}$". Conversely, the precision score is derived by comparing each token in "$\hat{x}$" to a token in "$x$". A greedy matching strategy is employed to pair tokens that demonstrate the highest degree of similarity. This method provides a comprehensive analysis of how precise and accurate the language model's output aligns with the factual reference sentence.

\paragraph{BLEURT}
the Bilingual Evaluation Understudy with Representations from Transformers (BLEURT) metric \citep{sellam2020bleurt} employs a unique pre-training scheme, where BERT is initially trained on a significant corpus of synthetic sentence pairs. This is reinforced by multiple lexical and semantic-level supervision signals, used concurrently. Following this pre-training stage, BERT is further fine-tuned on rating data, and its objective is to estimate human rating scores accurately. The initial stage of pre-training is essential for this metric because it significantly enhances the model's robustness, thereby effectively transforming it into a secure bulwark against quality drifts inherent in generative systems.

\paragraph{BARTScore} ~\citep{yuan2021bartscore} is a metric proposed to evaluate the quality of the generated text, such as in applications like machine translation and summarization. This metric conceives this evaluation as a problem of text generation, modeled using pre-trained sequence-to-sequence models. BARTScore uses BART, an encoder-decoder-based pre-trained model, to translate generated text to and from a reference point, earning higher scores when the text is more accurate and fluent. This metric offers several variants that can be flexibly applied in an unsupervised manner for different perspectives of text evaluation, such as fluency or factuality. Tests have shown that BARTScore can outperform existing top-scoring metrics in the majority of test settings across multiple datasets and perspectives. 


\subsubsection{Human Evaluation Metrics}

Human evaluation in factuality assessment is crucial due to its sensitivity to nuanced elements of language and context that may elude automated systems. Human evaluators excel at interpreting abstract concepts and emotional subtleties that can significantly inform the accuracy of evaluation. However, they are subject to limitations such as subjectivity, inconsistency, and potential for error. On the other hand, automated evaluations offer consistent results, and efficient processing of large data sets, and are ideal for tasks needing quantitative measurements. They also provide an objective benchmark for model performance comparison. Overall, an ideal evaluation framework might blend automated evaluation's scalability and consistency with human evaluation's ability to interpret complex linguistic concepts.


\paragraph{Attribution} is a metric to verify that the output of LLMs is sharing only verifiable information about the external world. 
As proposed by \citep{rashkin2023measuring}, Attributable to Identified Sources (AIS) is a human evaluation framework that adopts a binary concept of attribution. A text passage $y$ is deemed attributable to a set $A$ of evidence if, and only if, an arbitrary listener would agree with the statement "According to $A$, $y$" within the context of $y$. The AIS framework awards a full score (1.0) if every element of content in passage $y$ can be linked to the evidence set $A$. Conversely, it gives a score of zero (0.0) if this condition is not met.

Based on AIS,~\citet{gao2023rarr} propose a more fine-grained, sentence-level extension of AIS called Auto-AIS, where annotators assign AIS scores to each sentence, and an average score across all sentences is reported. This procedure effectively measures the percentage of sentences that are fully attributable to the evidence. Context, such as surrounding sentences and the question the text answers, is provided to annotators for more informed judgment. A limit is also set for the number of evidence snippets in the attribution report to maintain conciseness.

During model development, an automated AIS metric is defined to approximate human AIS evaluations, using a natural language inference model, which correlates well with AIS scores. Before computing the scores, they improve accuracy by decontextualizing each sentence in context. 

\paragraph{FActScore} \citep{Min2023-FActScore} is a novel evaluation metric designed to assess the factual precision of long-form text generated by LLMs. The challenge of evaluating the factuality of such text arises from two main issues: (1) the generated content often contains a mix of supported and unsupported information, making binary judgments insufficient, and (2) human evaluation is both time-consuming and expensive.
To address these challenges, FActScore breaks down a generated text into a series of atomic facts—short statements that each convey a single piece of information. Each atomic fact is then evaluated based on its support from a reliable knowledge source. The overall score represents the percentage of atomic facts that are supported by the knowledge source.
The paper conducted extensive human evaluations to compute FActScores for biographies generated by several state-of-the-art commercial LLMs, including InstructGPT, ChatGPT, and the retrieval-augmented PerplexityAI. The results revealed that these LLMs often contain factual inaccuracies, with FActScores ranging from 42\% to 71\%. Notably, the factual precision of these models tends to decrease as the rarity of the entities in the biographies increases.




\subsubsection{LLM-based Metrics}
Using LLMs for evaluation offers efficiency, versatility, reduced reliance on human annotation, and the capability of evaluating multiple dimensions of conversation quality in a single model call, which improves scalability. However, potential issues include a lack of established validation, which can lead to bias or accuracy problems if the LLM used for evaluation is not thoroughly vetted. The decision process to identify suitable LLMs and decoding strategies can be complex and pivotal to obtaining accurate evaluations. The range of evaluation may also be limited, as the focus is often on open-domain conversations, possibly leaving out assessments in specific or narrow domains. While reducing human input can be beneficial, it can also miss out on crucial interaction quality aspects better evaluated by human judges, such as emotional resonance or nuanced understandings.


\paragraph{GPTScore}~\citep{fu2023gptscore} is a new evaluation framework designed to assess the quality of output from generative AI models. To provide these evaluations, GPTScore taps into the emergent capabilities of 19 different pre-trained models, such as zero-shot instruction, and uses them to judge the generated texts. These models vary in scale from 80M to 175B. Testing across four text generation tasks, 22 aspects of evaluation, and 37 related datasets, has demonstrated that GPTScore can effectively evaluate text per instructions in natural language. This attribute allows it to sidestep challenges traditionally encountered in text evaluation, like the need for sample annotations and achieving custom, multi-faceted evaluations.

\paragraph{GPT-judge}~\citep{TruthfulQA} is a finetuned model based on the GPT-3-6.7B, which is trained to evaluate the truthfulness of answers to questions in the TruthfulQA dataset. The training set consists of triples in the form of question-answer-label combinations, where the label can be either true or false. The model's training set includes examples from the benchmark and answers generated by other models assessed by human evaluation. In its final form, the GPT-judge uses examples from all models to evaluate the truthfulness of responses. This training includes all questions from the dataset, with the goal being to evaluate truth, not generalize new questions.


The study conducted by the authors~\citep{TruthfulQA} focuses on the application of GPT-judge in assessing \emph{Truthfulness} and \emph{Informativeness} using the TruthfulQA Dataset. The authors undertook the fine-tuning of two distinct GPT-3 models to evaluate two essential aspects: Truthfulness, which pertains to the accuracy and honesty of information provided by the LLM, and Informativeness, which measures how effectively the LLM conveys relevant and valuable information in its responses.
From these two fundamental concepts, the authors derived a combined metric denoted as \emph{truth * info}. This metric represents the product of scalar scores for both truthfulness and informativeness. It not only quantifies the extent to which questions are answered truthfully but also incorporates the assessment of informativeness for each response. This comprehensive approach prevents the model from generating generic responses like "I have no comment" and ensures that responses are not only truthful but also valuable.
These metrics have found widespread deployment in evaluating the factuality of information generated by LLMs~\citep{chuang2023dola, li2023inferencetime}.

\paragraph{LLM-Eval} \citep{lin2023llm} is a novel evaluation methodology for open-domain dialogues with LLMs. Unlike conventional evaluation methods which rely on human annotations, ground-truth responses, or multiple LLM prompts, LLM-Eval uses a unique prompt-based evaluation process employing a unified schema to assess various elements of a conversation's quality during a single model function. Extensive evaluations of LLM-Eval's performance using multiple benchmark datasets indicate it is effective, efficient, and adaptable compared to traditional evaluation practices. Further, the authors stress the necessity of selecting appropriate LLMs and decoding strategies for precise evaluation outcomes, underscoring LLM-Eval's versatility and dependability in assessing open-domain conversation systems across a variety of circumstances.





\begin{table*}[]\small
\caption{Benchmarks for Evaluating Factuality in LLMs. This table details each benchmark's task type, associated sub-datasets, evaluation metrics, and the performance of select representative LLMs.}% The column `Human Annotation' indicates the benchmark involves human annotation.}
\centering
\label{tab:eval_benchmark}
\vspace{-2mm}
\begin{adjustbox}{max width=\textwidth}
  \setlength{\tabcolsep}{2mm}
  \renewcommand{\arraystretch}{1.25}
{

\begin{tabular}{lllll}
\toprule
Reference & Task Type & Dataset & Metrics & \makecell[l]{Performance of\\ Representative LLMs}  \\
\midrule 
% 

MMLU \citep{MMLU}
&
\begin{tabular}[c]{@{}l@{}}
Multi-Choice QA
\end{tabular}&
\begin{tabular}[c]{@{}l@{}}
Humanities,\\
Social,\\Sciences,\\
STEM...
\end{tabular}&
ACC&
\begin{tabular}[c]{@{}l@{}}
(ACC, 5-shot)\\
GPT-4: 86.4\quad  GPT-3.5: 70  \\LLaMA2-70B: 68.9
\end{tabular} 
\\
\hline 

TruthfulQA \citep{TruthfulQA} &
\begin{tabular}[c]{@{}l@{}}QA \end{tabular}&
\begin{tabular}[c]{@{}l@{}}Health, Law,\\ Conspiracies,\\ Fiction...\end{tabular}&
\begin{tabular}[c]{@{}l@{}}Human Score,\\ GPT-judge, \\ ROUGE, BLEU, \\ MC1,MC2... \end{tabular}&
\begin{tabular}[c]{@{}l@{}}
(zero-shot)\\
GPT-4: $\sim$29 (MC1)\\ GPT-3.5: $\sim$28 (MC1),  79.92(\%true)  \\
LLaMA2-70B: 53.37 (\%true)
\end{tabular}
\\
\hline

C-Eval \citep{C-Eval} &
Multi-Choice QA &
\begin{tabular}[c]{@{}l@{}} 
STEM,\\Social Science,\\Humanities...
\end{tabular} &
\begin{tabular}[c]{@{}l@{}} 
ACC
\end{tabular} &
\begin{tabular}[c]{@{}l@{}} 
(zero-shot, average ACC)\\
GPT-4: 68.7\quad GPT-3.5: 54.4\\
LLaMA2-70B: 50.13
\end{tabular} 
 \\
 \hline


AGIEval \citep{C-Eval} &
Multi-Choice QA &
\begin{tabular}[c]{@{}l@{}} 
Gaokao, \\(geometry, Bio,\\ history...),\\SAT, Law...
\end{tabular} &
\begin{tabular}[c]{@{}l@{}} 
ACC
\end{tabular} &
\begin{tabular}[l]{@{}l@{}} 
(zero-shot, average ACC)\\
GPT-4: 56.4\quad GPT-3.5: 42.9\\
LLaMA2-70B: 40.02
\end{tabular} 
\\
\hline


HaluEval \citep{HaluEval} &
\begin{tabular}[c]{@{}l@{}} 
Hallucination\\Evaluation \end{tabular}&
HaluEval &
ACC&
\begin{tabular}[c]{@{}l@{}} 
(general ACC)\\
GPT-3.5: 86.22
\end{tabular}
\\
\hline

BigBench \citep{BigBench} &
\begin{tabular}[c]{@{}l@{}} 
Multi-tasks
(QA, NLI, \\Fact Checking,\\ Reasoning...)
\end{tabular}&
BigBench
&
\begin{tabular}[c]{@{}l@{}} 
Metric to each\\type of task
\end{tabular}&
\begin{tabular}[c]{@{}l@{}} 
(Big-Bench Hard)\\
GPT-3.5: 49.6 \\
LLaMA-65B: 42.6
\end{tabular}
\\
\hline


ALCE \citep{gao2023enabling} &
\begin{tabular}[c]{@{}l@{}} 
Citation\\Generation
\end{tabular} &
\begin{tabular}[c]{@{}l@{}} 
ASQA, ELI5,\\QAMPARI\\
\end{tabular} &
\begin{tabular}[c]{@{}l@{}} 
MAUVE,\\
Exact Match,\\
ROUGE-L...
\end{tabular} &
\begin{tabular}[c]{@{}l@{}} 
(ASQA, 3-psg, citation prec)\\
GPT-3.5: 73.9\\
LLaMA-33B: 23.0
\end{tabular} 
\\
\hline

QUIP \citep{weller2023according} &
Generative QA &
\begin{tabular}[c]{@{}l@{}}
 TriviaQA,\\ NQ, ELI5,\\HotpotQA 
 \end{tabular} &
\begin{tabular}[c]{@{}l@{}}
 QUIP-Score \\ Exact match
 \end{tabular} &
 \begin{tabular}[c]{@{}l@{}}
 (ELI5, QUIP, null prompt)\\
 GPT-4: 21.0\\ GPT-3.5: 27.7
 \end{tabular} 
 \\
  \hline

 PopQA \citep{PopQA} &
 Multi-Choice QA &
\begin{tabular}[c]{@{}l@{}}
PopQA,\\
EntityQuestions
 \end{tabular} &
 ACC &
\begin{tabular}[c]{@{}l@{}}
(overall ACC)\\
GPT-3.5: $\sim$37.0
\end{tabular} 
 \\
 \hline

UniLC \citep{UniLC} &
Fact Checking &
\begin{tabular}[c]{@{}l@{}}
Climate,\\
Health, MGFN
 \end{tabular} &
 \begin{tabular}[c]{@{}l@{}}
ACC\\F1
 \end{tabular} &
\begin{tabular}[c]{@{}l@{}}
(zero-shot, fact tasks, average F1)\\
GPT-3.5: 51.62
\end{tabular} 
\\
\hline

Pinocchio \citep{Pinocchio} &
\begin{tabular}[c]{@{}l@{}}
Fact Checking\\QA, Reasoning \end{tabular}&
\begin{tabular}[c]{@{}l@{}}
Pinocchio
 \end{tabular} &
 \begin{tabular}[c]{@{}l@{}}
ACC\\F1
 \end{tabular} &
\begin{tabular}[c]{@{}l@{}}
GPT-3.5: (Zero-shot ACC: 46.8, F1:44.4)\\
GPT-3.5: (Few-shot ACC: 47.1, F1:45.7)
\end{tabular} 
\\
\hline


SelfAware \citep{yin-etal-2023-large} &
 Self-evaluation&
\begin{tabular}[c]{@{}l@{}}
SelfAware
 \end{tabular} &
 \begin{tabular}[c]{@{}l@{}}
ACC
 \end{tabular} &
\begin{tabular}[c]{@{}l@{}}
(instruction input, F1)\\
GPT-4: 75.47\quad GPT-3.5: 51.43\\
LLaMA-65B: 46.89
\end{tabular}
\\
\hline

RealTimeQA \citep{kasai2022realtimeqa} &
\begin{tabular}[c]{@{}l@{}} 
Multi-Choice QA,\\Generative QA \end{tabular}&
RealTimeQA &
ACC, F1 &
\begin{tabular}[c]{@{}l@{}}
(original setting, GCS retrieval)\\
GPT-3: 69.3 (ACC for MC)\\
GPT-3: 39.4 (F1 for generation)
\end{tabular}
\\
 \hline

FreshQA \citep{vu2023freshllms} &
Generative QA&
FRESHQA&
ACC (Human)&
\begin{tabular}[c]{@{}l@{}}
(strict ACC, null prompt)\\
GPT-4: 28.6\quad GPT-3.5: 26.0
\end{tabular}
\\

\bottomrule
\end{tabular}}
\end{adjustbox}
\end{table*}

% \subsection{Evaluating Factuality of LLMs}

\subsection{{Benchmarks for Factuality Evaluation}}
\label{sec:eval_benchmark}

In this section, we delve into the benchmarks that are prominently employed to assess the factuality of LLMs.
Specific benchmarks tailored for evaluating factuality in LLMs are tabulated in Table \ref{tab:eval_benchmark}. 

\begin{table*}[]\small
\caption{Factuality Evaluation Studies of LLMs. This table shows studies focused on factuality evaluation, and those providing valuable insights into such evaluations. The column `Human Eval' denotes if the evaluation incorporated human evaluation. The `Granularity' column specifies the level of evaluation, with token-level (T) and sentence-level (S) distinctions.}
\centering
\label{tab:evaluation1}
\vspace{-2mm}
\begin{adjustbox}{max width=\textwidth}
  \setlength{\tabcolsep}{0.9mm}
  \renewcommand{\arraystretch}{1}
{

\begin{tabular}{llllclc}
\toprule
Reference & Task & Dataset & Metrics & \makecell[l]{Human\\Eval} & Evaluated LLMs & Granularity \\
\midrule 
\begin{tabular}[c]{@{}l@{}} 
FActScore\\
\citep{Min2023-FActScore} 
\end{tabular}&
\begin{tabular}[c]{@{}l@{}}Biography\\Generation \end{tabular}&
\begin{tabular}[c]{@{}l@{}}183 people\\entities \end{tabular}&
F1&
\checkmark 
&
\begin{tabular}[c]{@{}l@{}} 
GPT-3.5\\
ChatGPT...
\end{tabular}
&
T
\\
\hline

\begin{tabular}[c]{@{}l@{}} 
SelfCheckGPT\\ \citep{SelfCheckGPT}
\end{tabular}&
Bio Generation&
WikiBio&
\begin{tabular}[c]{@{}l@{}}
AUC-PR,\\
Human Score
\end{tabular}&
\checkmark
&
\begin{tabular}[c]{@{}l@{}}
GPT-3,\\
LLaMA,\\
OPT,\\
GPT-J...
\end{tabular}
&
S
\\
\hline

\citet{wang2023evaluating} &
Open QA &
\begin{tabular}[c]{@{}l@{}}NQ, TQ  \\ \end{tabular}&
\begin{tabular}[c]{@{}l@{}}ACC, \\ EM\end{tabular} &
\checkmark &
\begin{tabular}[c]{@{}l@{}}GPT-3.5\\ ChatGPT\\GPT-4\\Bing Chat\end{tabular}&
S 
\\
\hline

\citet{pezeshkpour2023measuring} &
\begin{tabular}[c]{@{}l@{}}Knowledge \\Probing \end{tabular}&
\begin{tabular}[c]{@{}l@{}}T-REx, \\LAMA\end{tabular}&
ACC &
&
GPT3.5 &
T 
\\
\hline
    % \citet{borgeaud2022improving} &
    % \begin{tabular}[c]{@{}l@{}} QA, \\Language\\Modeling\end{tabular} &
    % \begin{tabular}[c]{@{}l@{}}
    % MassiveText, \\
    % C4, \\
    % Curation Corpus, \\
    % Wikitext103, \\
    % Lambada, \\
    % Pile, \\
    % NaturalQA
    % \end{tabular} &
    % \begin{tabular}[c]{@{}l@{}} 
    % PPL, \\ACC, \\Exact Match
    % \end{tabular} &
    %    \checkmark&
    %    \begin{tabular}[c]{@{}l@{}} 
    % Retro
    % \end{tabular} &
    % T 
    % \\
\citet{de-cao-etal-2021-editing} &
\begin{tabular}[c]{@{}l@{}} 
QA, \\
Fact Checking \end{tabular}&
\begin{tabular}[c]{@{}l@{}} 
KILT, \\FEVER, \\zsRE
\end{tabular} &
ACC &
&
\begin{tabular}[c]{@{}l@{}} GPT-3\\
FLAN-T5\end{tabular} &
\begin{tabular}[c]{@{}l@{}} S/
T \end{tabular} 
\\
\hline

 \citet{varshney2023stitch}&
 \begin{tabular}[c]{@{}l@{}} Article\\Generation\end{tabular}&
 \begin{tabular}[c]{@{}l@{}} Unnamed\\Dataset\end{tabular}&
 \begin{tabular}[c]{@{}l@{}} 
 ACC, \\AUC
 \end{tabular}&
 &
 \begin{tabular}[c]{@{}l@{}} 
 GPT3.5\\
 Vicuna
 \end{tabular} &
 S 
 \\
 \hline

\begin{tabular}[c]{@{}l@{}} 
FactTool\\
\citep{chern2023factool}
\end{tabular}&
KB-based QA&
\begin{tabular}[c]{@{}l@{}} 
RoSE
\end{tabular}&
\begin{tabular}[c]{@{}l@{}}
ACC\\F1...
\end{tabular}&
&
\begin{tabular}[c]{@{}l@{}} 
GPT-4\\
ChatGPT \\
FLAN-T5
\end{tabular}&
S
\\
\hline

\citet{kadavath2022language} &
Self-evaluation &
\begin{tabular}[c]{@{}l@{}} 
BIG Bench,\\
MMLU, 
LogiQA,\\
TruthfulQA,\\
QuALITY\\
TriviaQA\\
Lambada
\end{tabular}&
\begin{tabular}[c]{@{}l@{}} 
ACC, \\
Brier Score, \\
RMS\\Calibration\\Error...
\end{tabular}&
&
Claude&
T
\\ \midrule
\citet{chen2023beyond} &
QA &
\begin{tabular}[c]{@{}l@{}} 
NQ,\\
WoW
\end{tabular}&
\begin{tabular}[c]{@{}l@{}} 
Factuality, \\Relevance,\\ Coherence,\\ Informativeness,\\ Helpfulness,\\ and Validity
\end{tabular}&
&
LLaMA, FLAN-T5, ...&
T
\\
\midrule
RAGAS \citep{es2023ragas} &
QA &
\begin{tabular}[c]{@{}l@{}} 
WikiEval
\end{tabular}&
\begin{tabular}[c]{@{}l@{}} 
Faithfulness, \\Answer Relevance,\\ Context Relevance
\end{tabular}&
&
ChatGPT&
T
\\
\bottomrule
\end{tabular}}
\end{adjustbox}
\end{table*}


MMLU \citep{MMLU} and TruthfulQA \citep{TruthfulQA} stand as two pivotal benchmarks in the realm of evaluating the factuality of LLMs \citep{GPT4,llama2,Open-LLM-Leaderboard} . 
The MMLU benchmark is proposed to measure a text model's multitask accuracy across a diverse set of 57 tasks. These tasks span a wide range of subjects, from elementary mathematics to US history, computer science, law, and more. The benchmark is designed to test both the world knowledge and problem-solving ability of models. The findings from the paper suggest that while most recent models perform at near random-chance accuracy, the largest GPT-3 model showcased a significant improvement. However, even the best models still have a long way to go before achieving expert-level accuracy across all tasks \citep{GPT4}.
TruthfulQA is a benchmark designed to assess the truthfulness of a language model's generated answers. The benchmark consists of 817 questions spanning 38 categories, including health, law, finance, and politics. These questions were crafted in such a way that some humans might answer them falsely due to misconceptions or false beliefs. The goal for models is to avoid generating these false answers that they might have learned from imitating human texts. The TruthfulQA benchmark serves as a tool to highlight the potential pitfalls of relying solely on LLMs for accurate information and emphasizes the need for continued research in this area.


HaluEval \citep{HaluEval} is a benchmark designed to understand and evaluate the propensity of LLMs like ChatGPT to generate hallucinations. A hallucination, in this context, refers to content that either conflicts with the source or cannot be verified based on factual knowledge.
The HaluEval benchmark offers a vast collection of generated and human-annotated hallucinated samples, aiming to evaluate the performance of LLMs in recognizing such hallucinations. The benchmark utilizes a two-step framework, termed "sampling-then-filtering", based on ChatGPT to generate these samples. Additionally, human labelers were employed to annotate hallucinations in ChatGPT responses. The HaluEval benchmark is a comprehensive tool that not only evaluates the hallucination tendencies of LLMs but also provides insights into the types of content and the extent to which these models are prone to hallucinate. 

BigBench \citep{BigBench} focuses on the capabilities and limitations of LLMs. It comprises 204 tasks from diverse domains such as linguistics, childhood development, math, common-sense reasoning, biology, physics, social bias, software development, and more. The benchmark is designed to evaluate tasks believed to be beyond the capabilities of current language models. The study evaluates the performance of various models, including OpenAI's GPT models, on BIG-bench and compares them to human expert raters. Key findings suggest that model performance and calibration improve with scale but are still suboptimal when compared to human performance. Tasks that involve a significant knowledge or memorization component show predictable improvement, while tasks that exhibit "breakthrough" behavior at a certain scale often involve multiple steps or components. 

\citet{C-Eval} propose C-Eval, the first comprehensive Chinese evaluation suite. It can be used to evaluate the advanced knowledge and reasoning abilities of foundational models within a Chinese context. The evaluation suite comprises multiple-choice questions spanning 52 diverse disciplines, with four levels of difficulty: middle school, high school, college, and professional. Additionally, C-Eval Hard is introduced for very challenging subjects within the C-Eval suite, which demands advanced reasoning abilities to solve. Their evaluations of state-of-the-art LLMs, including both English and Chinese-oriented models, have shows that there is significant room for improvement as only GPT-4 managed to achieve an average accuracy of over 60\%.
The authors focus on assessing LLMs' advanced abilities within a Chinese context. The researchers assert that LLMs intended for a Chinese environment should be evaluated based on their knowledge of Chinese users' primary interests, such as Chinese culture, history, and laws.
With C-Eval, the authors aim to guide developers in understanding the abilities of their models from multiple dimensions to facilitate the development and growth of foundational models for Chinese users. Simultaneously, C-Eval has not just introduced a whole suite but subsets that can serve as individual benchmarks, thereby assessing certain model abilities and analyzing key strengths and limitations of foundational models. Experiment results have shown that although GPT-4, ChatGPT, and Claude were not exclusively tailored for Chinese data, they emerged as the top performers on C-Eval.

SelfAware \citep{yin-etal-2023-large} aims to investigate whether models recognize what they don't know. This dataset encompasses two types of questions: unanswerable and answerable. The dataset comprises 2,858 unanswerable questions gathered from various websites and 2,337 answerable questions extracted from sources such as SQuAD, HotpotQA, and TriviaQA. Each unanswerable question is confirmed as such by three human evaluators. In the conducted experiments, GPT-4 achieves the highest F1 score of 75.5, compared to a human score of 85.0. Larger models tend to perform better and in-context learning can enhance performance.

The Pinocchio benchmark \citep{Pinocchio} serves as an extensive evaluation platform, emphasizing factuality and reasoning for LLMs. This benchmark encompasses 20,000 varied factual queries from diverse sources, timeframes, fields, regions, and languages. It tests an LLM's capability to discern combined facts, process both organized and scattered evidence, recognize the temporal evolution of facts, pinpoint minute factual disparities, and withstand adversarial inputs. Each reasoning challenge within the benchmark is calibrated for difficulty to allow for detailed analysis.


\citet{kasai2022realtimeqa} introduce a dynamic QA platform named REALTIMEQA. This platform is unique in that it announces questions and evaluates systems on a regular basis, specifically weekly. The questions posed by REALTIMEQA pertain to current events or novel information, challenging the static nature of traditional open domain QA datasets. The platform aims to address instantaneous information needs, pushing QA systems to provide answers about recent events or developments. Their preliminary findings indicate that while GPT-3 can often update its generation results based on newly-retrieved documents, it sometimes returns outdated answers when the retrieved documents lack sufficient information. 


FreshQA \citep{vu2023freshllms} is a dynamic benchmark designed to evaluate up-to-date world knowledge of LLMs. Its questions range from those that are never-changing to those that are fast-changing, as well as questions based on false premises. The aim is to challenge the static nature of LLMs and test their adaptability to ever-changing knowledge through human evaluations. They develop a reliable evaluation protocol that uses a two-mode system: RELAXED and STRICT for a comprehensive understanding of model performance, ensuring that answers are confident, definitive, and accurate. They also provide a strong baseline named FRESHPROMPT, which seeks to enhance LLM performance by integrating real-time data from search engines. Initial experiments reveal that, outdated training data weakens the performance of LLMs and the FRESHPROMPT method can significantly enhance it. The research underscores the need for LLMs to be refreshed with current information to ensure their relevance and accuracy in a constantly evolving world.

Several benchmarks, such as BigBench \citep{BigBench} and C-Eval \citep{C-Eval}, encompass subsets that extend beyond the realm of factual knowledge or factuality. In this work, we specifically emphasize and focus on those subsets related to factuality.

There are benchmarks primarily designed for Pre-trained Language Models (PLMs) that can also be adapted for LLMs. Some studies use them for evaluating LLM's factuality, but they are not so widely-used, so we have chosen to exclude them from the table for clarity. These benchmarks predominantly encompass knowledge-intensive tasks, as highlighted by \citep{kilt}. Those benchmarks include NaturalQuestions (NQ) \citep{NQ}, TriviaQA (TQ) \citep{TQ}, OTT-QA~\citep{OTT-QA}, AmbigQA~\cite{ambigQA} and WebQuestion (WQ) \citep{WQ} of Open-domain Question Answering (QA) task, the HotpotQA~\citep{hotpotqa}, 2WikiMultihopQA~\citep{2WikiMultihopQA}, IIRC~\citep{IIRC}, MuSiQue~\citep{musique} of multi-step
QA, the ELI5 \citep{ELI5} of the Long-form QA task, the FEVER \citep{fever}, FM2 \citep{FM2}, HOVER \citep{HOVER}, FEVEROUS \citep{FEVEROUS} of the Fact Checking task, the T-REx \citep{TREx}, zsRE \citep{ZsRE} and LAMA \citep{LAMA} to examine factual knowledge contained in pretrained language models, the WikiBio \citep{wikiBio} of biography generation task, the RoSE \citep{Rose}, WikiAsp \citep{wikiasp} of summarization task, the KILT \citep{kilt} of comprehensive knowledge intensive tasks, the MassiveText \citep{Gopher}, Curation Corpus \citep{curation}, Wikitext103 \citep{Wikitext103}, Lambada \citep{lambada}, C4 \citep{C4}, Pile \citep{pile} of langauge modeling task, the WoW \citep{WoW}, DSTC7 track2~\citep{dstc7}, DSTC11 track5~\citep{dstc11} of dialogue task, the RealToxicityPrompts \cite{RealToxicityPrompts} of toxicity reduction, the CommaQA \citep{CommaQA}, StrategyQA \citep{StrategyQA}, TempQuestions \citep{TempQuestions}, IN-FOTABS~\citep{INFOTABS} of diverse reasoning tasks.

%\textcolor{red}{TBC} 

Some studies \citep{Min2023-FActScore, SelfCheckGPT} also provide a small dataset, but they mainly concentrate on the evaluation metrics or methods for factuality, we choose to discuss them in the next subsection.

\subsection{Factuality Evaluation Studies}
\label{sec:factuality-eval}


%%%%%%%%%%%%%%%%%%%%%%%%%%%%%%%%%%%%%%%%%%
\begin{table*}[]\small
\caption{Factuality Evaluation of LLMs. This table lists the benchmarks that used in the enhancement work. The column `Human Eval' denotes if the evaluation incorporated human evaluation. The `Granularity' column specifies the level of evaluation, with token-level (T) and sentence-level (S) distinctions.}
\centering
\label{tab:eval-table2}
\vspace{-2mm}
\begin{adjustbox}{max width=\textwidth}
  \setlength{\tabcolsep}{1.2mm}
    \renewcommand{\arraystretch}{0.88}
{

\begin{tabular}{llllclc}
\toprule
Reference & Task & Dataset & Metrics & \makecell[l]{Human\\Eval} & Evaluated LLMs & Granularity \\
\midrule 
\begin{tabular}[c]{@{}l@{}}
Retro\\
\citep{borgeaud2022improving} 
\end{tabular} &
\begin{tabular}[c]{@{}l@{}} QA, \\Language\\Modeling\end{tabular} &
\begin{tabular}[c]{@{}l@{}}
MassiveText, \\
Curation Corpus, \\
Wikitext103, \\
Lambada, \\
C4,Pile, NQ
\end{tabular} &
\begin{tabular}[c]{@{}l@{}} 
PPL, \\ACC, \\Exact Match
\end{tabular} &
   \checkmark&
\begin{tabular}[c]{@{}l@{}} 
Retro
\end{tabular} &
T 
\\
\hline
\begin{tabular}[c]{@{}l@{}} 
GenRead\\
 \citep{yu2023generate} 
\end{tabular} &
\begin{tabular}[c]{@{}l@{}} QA, \\ Dialogue, \\ Fact Checking \end{tabular} &
%\makecell[l]{ NQ, TQ, WebQ, \\FEVER, \\FM2, WoW}&
\begin{tabular}[c]{@{}l@{}} NQ, TQ, WebQ, \\FEVER, \\FM2, WoW  \end{tabular} &
\begin{tabular}[c]{@{}l@{}} 
EM, ACC, \\F1, Rouge-L
\end{tabular} &
   &
   \begin{tabular}[c]{@{}l@{}} 
GPT3.5, Codex\\GPT-3, Gopher\\FLAN, GLaM\\PaLM\end{tabular} &
S 
\\
\hline
\begin{tabular}[c]{@{}l@{}}
GopherCite\\ \citep{menick2022teaching}  
\end{tabular}
&
\begin{tabular}[c]{@{}l@{}} Self-supported QA\end{tabular} &
\begin{tabular}[c]{@{}l@{}} 
NQ, ELI5, \\
 TruthfulQA\\
(Health\\Law, Fiction\\ Conspiracies)

\end{tabular} &
\begin{tabular}[c]{@{}l@{}} 
Human Score
\end{tabular} &
   \checkmark&
   \begin{tabular}[c]{@{}l@{}} 
GopherCite
\end{tabular} &
S 
\\
\hline
\citet{trivedi-etal-2023-interleaving} &
\begin{tabular}[c]{@{}l@{}} 
 QA
\end{tabular} &
\begin{tabular}[c]{@{}l@{}} 
HotpotQA, IIRC\\2WikiMultihopQA, \\MuSiQue(music)
\end{tabular} &
\begin{tabular}[c]{@{}l@{}} 
Retrieval recall, \\
Answer F1
\end{tabular} &
 &
 \begin{tabular}[c]{@{}l@{}} 
 GPT-3\\
FLAN-T5
 \end{tabular} &
  \begin{tabular}[c]{@{}l@{}} 
 S/
T
 \end{tabular} 
 \\
\hline
 \citet{peng2023check} &
\begin{tabular}[c]{@{}l@{}} 
QA, \\
Dialogue
\end{tabular} &
\begin{tabular}[c]{@{}l@{}} 
DSTC7 track2, \\
DSTC11 track5, \\
OTT-QA
\end{tabular} &
\begin{tabular}[c]{@{}l@{}} 
ROUGE, chrF, \\
BERTScore, \\
Usefulness, \\
Humanness...
\end{tabular} &
\checkmark &
ChatGPT &
  \begin{tabular}[c]{@{}l@{}} 
 S/
T
 \end{tabular} 
 \\
\hline
\begin{tabular}[c]{@{}l@{}}
CRITIC\\
\citep{gou2023critic}
\end{tabular} &
\begin{tabular}[c]{@{}l@{}}
QA,\\Toxicity Reduction
 \end{tabular} &
\begin{tabular}[c]{@{}l@{}}
 AmbigNQ, \\TriviaQA, \\HotpotQA, 
\\RealToxicityPrompts
 \end{tabular}  &
 \begin{tabular}[c]{@{}l@{}}
Exact Match, \\
maximum toxicity, \\
perplexity, \\
n-gram diversity, \\
AUROC...,
\end{tabular} &
&
\begin{tabular}[c]{@{}l@{}}
GPT-3.5\\
ChatGPT
\end{tabular} &
T 
\\
\hline
\citet{khot2023decomposed} &
\begin{tabular}[c]{@{}l@{}}
QA, \\
long-context\\ QA
\end{tabular} &
\begin{tabular}[c]{@{}l@{}}
CommaQA-E, \\
2WikiMultihopQA, \\ MuSiQue, \\HotpotQA
\end{tabular} &
\begin{tabular}[c]{@{}l@{}}
Exact Match, \\
Answer F1
\end{tabular} &
&
\begin{tabular}[c]{@{}l@{}}
GPT-3\\
FLAN-T5
\end{tabular} &
T 
\\
\hline
\begin{tabular}[c]{@{}l@{}} 
ReAct\\
\citep{yao2023react} 
\end{tabular} &
\begin{tabular}[c]{@{}l@{}} 
QA,\\ Fact Verification 
\end{tabular} &
\begin{tabular}[c]{@{}l@{}}
HotpotQA, \\FEVER
\end{tabular} &
\begin{tabular}[c]{@{}l@{}}
Exact Match, \\
ACC
\end{tabular} &
 &
\begin{tabular}[c]{@{}l@{}}
PaLM\\
GPT-3
\end{tabular} &
\begin{tabular}[c]{@{}l@{}}
S/T
\end{tabular}
\\
\hline
\citet{jiang2023active} &
\begin{tabular}[c]{@{}l@{}}
QA, \\Commonsense\\Reasoning, \\
long-form QA...
\end{tabular} &
\begin{tabular}[c]{@{}l@{}} 
2WikiMultihopQA, \\StrategyQA, 
\\ASQA, WikiAsp
\end{tabular} &
\begin{tabular}[c]{@{}l@{}} 
Exact Match, \\
Disambig-F1, \\ROUGE, \\
entity F1...
\end{tabular} &
&
GPT-3.5 &
T 
\\
\hline
\citet{lee2022factuality} &
\begin{tabular}[c]{@{}l@{}}
Open-ended\\Generation 
\end{tabular}&
FEVER &
\begin{tabular}[c]{@{}l@{}} 
Entity score, \\
Entailment\\Ratio, ppl...
\end{tabular}&
 &
 Megatron-LM &
 T
\\
\hline
\begin{tabular}[c]{@{}l@{}}
SAIL\\
\citep{luo2023sail} 
\end{tabular}&
\begin{tabular}[c]{@{}l@{}}
QA,\\
Fact Checking 
\end{tabular}&
UniLC&
\begin{tabular}[c]{@{}l@{}} 
ACC\\ F1
\end{tabular}&
&
\begin{tabular}[c]{@{}l@{}}
LLaMA
Vicuna\\
SAIL
\end{tabular}&
T
\\
\hline
\citet{he2022rethinking} &
\begin{tabular}[c]{@{}l@{}} 
Commonsense\\ Reasoning, \\Temporal\\ Reasoning, \\
Tabular\\Reasoning
\end{tabular} &
\begin{tabular}[c]{@{}l@{}}
StrategyQA, \\ TempQuestions, \\ IN-FOTABS
\end{tabular}&
ACC&
&
GPT-3 &
T
\\
\hline
\citet{pan-etal-2023-fact} &
Fact Checking&
\begin{tabular}[c]{@{}l@{}}
HOVER \\ FEVEROUS-S
\end{tabular}&
Macro-F1 &
&
\begin{tabular}[c]{@{}l@{}}
Codex\\
FLAN-T5
\end{tabular}&
S
\\
\hline
\citet{multiagent_debate}&
\begin{tabular}[c]{@{}l@{}} 
Biography,\\
MMLU
\end{tabular}&
\begin{tabular}[c]{@{}l@{}} 
Unnamed\\Biography\\Dataset, \\
MMLU
\end{tabular}&
\begin{tabular}[c]{@{}l@{}} 
ChatGPT\\Evaluator, \\ACC
\end{tabular}&
&
\begin{tabular}[c]{@{}l@{}} 
Bard\\ChatGPT
\end{tabular}&
S
\\
\hline 
\citet{asai2023selfrag} &
\begin{tabular}[c]{@{}l@{}} 
QA,\\
CItation 
Generation
\end{tabular}&
\begin{tabular}[c]{@{}l@{}} 
PopQA,\\Pubhealth,\\ASQA,\\ARC,\\Unamed\\Biography
\end{tabular}&
\begin{tabular}[c]{@{}l@{}} 
ACC,\\Factscore,\\Exact Match,\\
ROUGE,\\MAUVE...
\end{tabular}&
&
\begin{tabular}[c]{@{}l@{}} 
LLaMA2\\ChatGPT\\Perplexity.ai\\Alpaca\\SAIL
\end{tabular}&
S/T
\\
    
\bottomrule
\end{tabular}}
\end{adjustbox}
\end{table*}

In this section, we delve into studies that evaluate factuality in LLMs without introducing a specific benchmark, focusing primarily on those whose main contribution lies in the evaluation methodology. We spotlight works that have pioneered evaluation techniques, metrics, or have offered distinctive insights into the factuality evaluation of LLMs.



\citet{SelfCheckGPT} use an evaluation process that encompasses several key steps. Initially, synthetic Wikipedia articles are generated using GPT-3, focusing on individuals from the Wikibio dataset. Subsequently, manual annotation is performed at the sentence level, classifying sentences as ``Major Inaccurate" , ``Minor Inaccurate", or ``Accurate", with ``Major Inaccurate" denoting sentences unrelated to the topic. Passage-level scores are derived by averaging sentence-level labels, and identifying cases of total inaccuracies through score distribution analysis. Inter-annotator agreement is assessed using Cohen's $\kappa$ ~\cite{cohen1960coefficient}. Evaluation metrics primarily employ precision-recall curves (PR-Curves), distinguishing between ``Non-Factual Sentences," ``Non-Factual* Sentences" (a specific subset), and ``Factual Sentences." These PR-Curves elucidate the trade-off between precision and recall for different detection methods.

BSChecker \citep{BSChecker} is a factual errors detection framework that detects inaccuracies on triplet-level granularity. Its pipeline comprises a claim extractor and a checker. The claim extractor extracts triplets from model’s response as claims formed with (subject, predicate, object). Meanwhile, the checker compare these claims triplets with references, categorizing each triplet to one of three hallucination labels: Entailment, Contradiction or Neutral, drawing inspiration from the Natural Language Inference task. BSChecker offers a benchmark that distinguishes three task settings: Zero Context, Noisy Context and Accurate Context covering tasks of closed-book QA, retrieval-augmented generation, summarization, closed-QA and information extraction.

\citet{pezeshkpour2023measuring} propose a new metric for measuring whether a certain type of knowledge is present in LLM. This metric is based on information theory and measures knowledge by analyzing the probability distribution of predictions made by LLM before and after injecting target knowledge. The accuracy of GPT-3.5 in the Knowledge Probing task is tested on the T-REx \citep{TREx} and LAMA \citep{LAMA} datasets. %They find that \textcolor{red}{XXXXX}.


% \citet{de-cao-etal-2021-editing} 
% propose a framework to evaluate the performance of LLM editing, represented by the edited parameters $\Theta_0$, with several key metrics on two different tasks: closed-book fact-checking (FC) and closed-book question answering (QA).
% The metrics are:
% \textit{(1) Success Rate: }This measures how well the model successfully updates the knowledge in parameters; 
% \textit{(2) Retain Accuracy:} This is a measure of how well parameters retain the original predictions of the model; 
% \textit{(3) Equivalence Accuracy:} This measures how consistent the revised model parameters's predictions are for semantically equivalent inputs; and
% \textit{(4) Performance Deterioration:} This is a measure of how much the updated model's test performance deteriorates.
% The authors claim that, prior studies like \cite{sinitsin2019editable} only measured performance deterioration and success rate, leaving out retained accuracy and equivalence accuracy. \citet{zhu2020modifying} also fell short in measuring retained accuracy while only exploring accuracy against old or new facts. The authors emphasize the importance of retaining accuracy in production settings where model predictions have potentially been rigorously analyzed and tested. They aim to maintain all original predictions while effectively avoiding the need to recalibrate the models, which is crucial in scenarios where probability estimates are used downstream.

\citet{varshney2023stitch} point out a common real-world occurrence where users often ask questions based on false premises. These questions are challenging for state-of-the-art models.
This necessitated the creation of a new evaluation dataset. To this end, the authors have conducted a case study and compiled a set of 50 such adversarial questions, all of which the GPT-3.5 model answered incorrectly. The aim is to create a challenging experimental setup to assess the performance of models faced with such questions.
In order to enhance the evaluation, corresponding true premise questions were created for each of the false premise questions. This allowe for a holistic evaluation of the model's performances, taking into consideration both correct and incorrect premises.
The authors make sure to evaluate the complete answers given by the model for correct and incorrect questions - in this context, it's not enough for an answer to be partially correct, the entire answer needs to be accurate to be marked as correct. For example: For the false premise question ``Why does Helium have an atomic number of 1?", the corresponding true premise question is ``Why does Hydrogen have an atomic number of 1?".


FACTOOL \citep{chern2023factool} is a tool designed to function as a factuality detector, with the primary purpose of auditing generative chatbots and assessing the reliability of their outputs. This tool is employed to evaluate several contemporary chatbots, including GPT-4, ChatGPT, Claude, Bard \citep{bard}, and Vicuna \citep{vicuna2023}. Notably, FACTOOL itself leverages the capabilities of GPT-4.
For the evaluation process, the researchers have curated a diverse set of prompts: 30 from knowledge-based question answering (KB-QA), 10 each from code, math, and scientific domains. The KB-QA prompts were sourced from a prior study, code prompts were taken from HumanEval, math prompts from another distinct study, while the scientific prompts were crafted by the authors themselves.
The evaluation metrics included both claim-level and response-level accuracies for each chatbot. To offer a more comprehensive and equitable evaluation, a weighted claim-level accuracy is used. The weighting is determined based on the proportion of prompts from each category.
The findings are illuminating. GPT-4 emerge as the top performer in terms of both weighted claim-level factual accuracy and response-level accuracy among all the chatbots assessed. Another intriguing observation is that chatbots that underwent supervised fine-tuning, such as Vicuna-13B, exhibited commendable performance in standard scenarios like KB-QA. However, their performance dip in more intricate scenarios, including those involving math, code, and scientific queries.

\citet{wang2023evaluating} ask several LLMs, including ChatGPT, GPT-4 \citep{GPT4}, BingChat \citep{BingChat} to answer open questions from NaturalQuestions \citep{NQ} and TriviaQA \citep{TQ}. They manually estimate the accuracy of those LLMs on open question answering, and find that though LLMs can achieve nice performance but still far away from perfect. Besides, they evaluate whether the GPT-3.5 can assess the correctness of LLM-generated responses, and find negative results, even if the golden answer is also presented.
Similarity, \citet{fu2023large} ask LLMs, such as GPT-2 and GPT-4, to directly score the factuality of a summary, and find  no significant correlation  between LLM's factuality indicators and human evaluations.


\citet{kadavath2022language} investigate whether language models can evaluate the accuracy of their own assertions and predict which questions they can answer correctly. It is found that larger models are well-calibrated on diverse multiple-choice and true/false questions if given in the appropriate format.
The approach to self-evaluation on open-ended tasks is to ask the models to initially suggest answers, and then evaluate the probability (P[True]) that their answers are correct. This resulted in compelling performance, calibration, and scaling on a diverse range of tasks. Furthermore, self-evaluation performance improved when the models are allowed to consider many of their own suggestions before predicting the validity of a specific one.


\citet{yu2023generate} explore whether the internal knowledge of LLMs can replace the retrieved documents on knowledge intensive tasks. They ask LLMs, such as InstructGPT \citep{gpt35}, to directly generate contexts given a question rather than retrieving from database. They find the generated documents contain the golden answers more often than the top retrieved documents. Then they feed the generated docs and retrieved docs to the Fusion-in-Decoder model \citep{Fusion-in-Decoder} for knowledge-intensive tasks such as Open-domain QA \citep{NQ} and find the generated docs are more effective than the retrieved docs, suggesting that the LLMs contain enough knowledge for knowledge-intensive tasks.

\citet{menick2022teaching} propose a task named \textit{Self-supported QA} to evaluate LLMs' ability in also producing citations when generating answers. Authors ask humans to evaluate whether the responses of their proposed model GopherCite are plausible and whether they are supported by the accompanying quote evidence on datasets such as NQ, ELI5, TruthfulQA.



CONNER \citep{chen2023beyond}, a framework that evaluates LLMs as generators of knowledge. It focuses on six areas: Factuality, Relevance, Coherence, Informativeness, Helpfulness, and Validity. It evaluates whether the generated information can be backed by external proof (Factuality), is relevant to the user's query (Relevance), and is logically consistent (Coherence). It also checks if the knowledge provided is novel or surprising (Informativeness). The Extrinsic evaluation measures whether the knowledge enhances downstream tasks (Helpfulness) and its results are factually accurate (Validity). 

In the realm of factuality evaluation, the Model Editing task holds a unique position, focusing on refining the internal knowledge of models. This task comes with its own set of specialized evaluation metrics. Predominantly, the Zero-Shot Relation Extraction (zsRE) \cite{ZsRE} and CounterFact \cite{meng2022locating} serve as the primary benchmarks for assessing Model Editing techniques. When evaluating these methods, several key criteria emerge:
Reliability: Post-editing, the model should consistently generate the intended output.
Generalization: The model should be adept at producing the target output even when presented with paraphrased inputs.
Locality: Edits should be localized, ensuring that facts not related to the specific edit remain intact.
However, given the intricate web of interconnected facts, recent studies \citet{zhong2023mquake,yao2023editing,cohen2023evaluating} have advocated for a more holistic evaluation approach. They introduce broader criteria for fact updates, encompassing aspects like portability, logical generalization, among others.

RAGAS \citep{es2023ragas} proposes a framework for the reference-free evaluation of Retrieval-Augmented Generation (RAG) systems in Large Language Models (LLMs). This framework primarily assesses the RAG system's ability to identify relevant and key context paragraphs, the LLM's fidelity in utilizing these paragraphs, and the overall quality of the generated content. The framework focuses on three aspects of quality:
Faithfulness: The generated answers should be based on the given context. To evaluate faithfulness, the generated content is broken down into sentences, which are then transformed into one or more short assertions. Each assertion's faithfulness is assessed against the retrieved content, with the faithfulness score being the ratio of faithful assertions to the total number of assertions.
Answer Relevance: The generated answers should appropriately address the posed questions. This is evaluated by having the LLM generate potential questions from each piece of generated content. An embedding model then calculates the similarity between each potential question and the original question. The answer relevance score is the average of these similarities.
Context Relevance: The retrieved context should be highly relevant, containing as little irrelevant information as possible. This is assessed by using the LLM to extract sentences from the retrieved content that are relevant to the question. The context relevance score is then calculated based on the proportion of these relevant sentences to the total number of sentences retrieved.
The authors created a WikiEval dataset, consisting of question-context-answer triplets with human judgments. This allows for the measurement of how closely different evaluation frameworks align with human assessments of faithfulness, answer relevance, and context relevance. The paper compares two baselines: GPTScore, which requires ChatGPT to rate the three quality dimensions on a scale of 0 to 10, and GPT Ranking, which does not require ChatGPT to choose a preferred answer/context but includes definitions of the considered quality metrics in the prompt. When evaluating answer relevance, a specific prompt is used.
Overall, the RAGAS framework proposed in the paper aligns more closely with human judgment compared to the two baselines. 



%%%%%%%%%

Some work's main contributions lie at the methods to improve the factuality of LLMs, and their evaluation part may also be informative and important when people reach the related study. So we choose to list evaluation parts in the Table \ref{tab:eval-table2} but not to discuss their evaluation in detail.
% Those work include \citet{borgeaud2022improving}

%%%%%%%%%
\subsection{Evaluating Domain-specific Factuality}


%In our examination of specialized Large Language Models (LLMs), we've identified an extensive array of datasets and benchmarks tailored to a variety of domains. These resources are instrumental not only in assessing the performance of LLMs but also in propelling forward domain-specific applications. For a comprehensive overview, we've compiled these resources in Table~\ref{tab:domain-eval}. The primary focus of this subsection is on domain-specific factuality evaluation, setting it apart from Sec \ref{sec:eval_benchmark} and \ref{sec:factuality-eval}. While the latter sections address general factuality evaluation, they only partly touch upon datasets geared towards factuality within specific domains. %Detailed information about these datasets and benchmarks can be found in the Appendix~\ref{sec:domain-eval}, allowing the main content to remain succinct and focused.

\label{sec:domain-eval}
\begin{table*}[]\small
\centering
\caption{Benchmarks for domain-specific factuality evaluation. The table presents the domain, tasks, datasets, and the LLMs evaluated in the respective study.}
\vspace{-2mm}
\label{tab:domain-eval}
\begin{adjustbox}{max width=\textwidth}
  \setlength{\tabcolsep}{2mm}
{
\begin{tabular}{lllll}
\toprule
Reference &
  Domain &
  Task &
  Metrics &
  %Human &
  Evaluated LLMs \\
  \midrule FLARE \citep{xie2023pixiu} &
  Finance  &
  \begin{tabular}[c]{@{}l@{}}Sentiment analysis, \\ News headline classification \\ Named entity recognition \\ Question answering \\ Stock movement prediction\end{tabular} &
  \begin{tabular}[c]{@{}l@{}}F1, ACC,\\Avg F1, \\ Entity F1, \\ EM, MCC\end{tabular} &
  %\checkmark &
  \begin{tabular}[c]{@{}l@{}}GPT-4 , \\ BloombergGPT, \\  FinMA-(7B, 30B, 7B-full), \\ Vicuna-7B \end{tabular}\\
   \midrule EcomInstruct \citep{li2023ecomgpt} & Finance  
   & 134 E-com tasks
   & \begin{tabular}[c]{@{}l@{}} Micro-F1, \\Macro-F1, \\
  ROUGE\end{tabular}
   & \begin{tabular}[c]{@{}l@{}}BLOOM, BLOOMZ,\\ ChatGPT, EcomGPT\end{tabular}
   \\
   \midrule CMB \citep{wang2023cmb} &
  Medicine &
  Multi-Choice QA &
  ACC &
  \begin{tabular}[c]{@{}l@{}}GPT-4, ChatGLM2-6B,\\ ChatGPT, DoctorGLM,\\ Baichuan-13B-chat,\\ HuatuoGPT, MedicalGPT,\\ ChatMed-Consult,\\ ChatGLM-Med ,\\ Bentsao, BianQue-2\end{tabular} \\
   \midrule Huatuo-26M \citep{li2023huatuo} &
Medicine & Generative-QA
    &
  \begin{tabular}[c]{@{}l@{}}BLEU, \\ROUGE, \\ GLEU\end{tabular} &
  T5, GPT2 \\
   \midrule NCBI \citep{jin2023genegpt} &
  Medicine &
  \begin{tabular}[c]{@{}l@{}}Nomenclature, \\ Genomic location, \\
  Functional analysis, \\
  Sequence alignment \end{tabular} &
  ACC & 
   \begin{tabular}[c]{@{}l@{}}GPT-2, BioGPT, \\
   BioMedLM, 
   GPT-3, \\
   ChatGPT, New Bing
   \end{tabular} \\
   \midrule LegalBench \citep{guha2023legalbench} &
  Law &
  \begin{tabular}[c]{@{}l@{}}Issue-spotting,  \\
  Rule-recall, \\
  Rule-application, \\ Rule-conclusion,
  \\
  Interpretation,
  \\
  Rhetorical-understanding
  \end{tabular}  &
  \begin{tabular}[c]{@{}l@{}} 
          ACC, EM
   \end{tabular}
  & 
   \begin{tabular}[c]{@{}l@{}}
   GPT-4, GPT-3.5, \\
   Claude-1, Incite, OPT\\
   Falcon, LLaMA-2, FLAN-T5...
   \end{tabular} \\
   \midrule 
   LawBench \citep{fei2023lawbench} &
  Law &
  \begin{tabular}[c]{@{}l@{}}Legal QA, NER, \\
  Sentiment Analysis, \\
  Reading Comprehension
  \end{tabular}  &
  \begin{tabular}[c]{@{}l@{}} 
         F1, ACC,\\
        ROUGE-L,\\
         Normalized \\log-distance	...
   \end{tabular}
  & 
   \begin{tabular}[c]{@{}l@{}}GPT-4,\\ 
   ChatGPT, \\
   InternLM-Chat,\\
   StableBeluga2...
   \end{tabular} \\
   \midrule 
   OceanBench \citep{bi2023oceangpt} &
  Geoscience &
  \begin{tabular}[c]{@{}l@{}}15  ocean-\\related tasks
  \end{tabular}  &
  \begin{tabular}[c]{@{}l@{}} 
         ACC
   \end{tabular}
  & 
   \begin{tabular}[c]{@{}l@{}}GPT-4
   \end{tabular} \\
  \bottomrule
\end{tabular}
}
\end{adjustbox}

\end{table*}

In our examination of specialized Large Language Models (LLMs), we've identified an extensive array of datasets and benchmarks tailored to a variety of domains.  These resources not only serve as critical tools for evaluating the capabilities of LLMs but also facilitate advancements in specialized applications. We summarize them in Table~\ref{tab:domain-eval}. The distinction between this subsection and the previous two lies in its focus. This subsection delves deeper into factuality evaluation tailored to specific domains, while Sec \ref{sec:eval_benchmark} and \ref{sec:factuality-eval} primarily concentrate on general factuality evaluation, with only a portion of their content dedicated to datasets evaluating factuality within specific domains.


{\paragraph{Finance}}
\citet{xie2023pixiu} designed a financial natural language understanding and prediction evaluation benchmark dubbed FLARE, based on their collected financial instruction tuning dataset FIT. This benchmark is used to evaluate their FinMA model. It randomly selects validation sets from FIT to choose the best model checkpoint and utilizes distinct test sets for evaluation.
FLARE is a broader variant compared to the existing FLUE benchmark~\cite{shah2022flue} as it also encapsulates financial prediction tasks like stock movement prediction in addition to standard NLP tasks. 
The FLARE dataset includes several subtasks, such as sentiment analysis (FPB, FiQA-SA), news headline  classification (Headline), named entity recognition (NER), question answering (FinQA, ConvFinQA), and stock movement prediction (BigData22, ACL18, CIKM18).
Performance is gauged via a variety of metrics for each task, such as the accuracy and weighted F1 Score for sentiment analysis, entity-level F1 score for named entity recognition, and accuracy and Matthews correlation coefficient for stock movement prediction.
Several methods, including their own FIT-fine-tuned FinMA and other LLMs (BloombergGPT, GPT-4, ChatGPT, BLOOM, GPT-NeoX, OPT-66B, Vicuna-13B) are dedicated to their comparison. BloombergGPT's performance is assessed in various shot scenarios, while the zero-shot performance is reported for the remaining results. Some of the baselines depend on human evaluations given that LLMs without fine-tuning fail to generate instruction-defined answers. Conversely, FinMA's results are conducted on a zero-shot basis and can be evaluated automatically. 
To enable direct comparison between the performance of FinMA and BloombergGPT, despite the former not releasing their test datasets, test datasets were constructed with the same data distribution.

\citet{li2023ecomgpt} propose EcomInstruct benchmark, which experiment investigates the performance of their EcomGPT language model in comparison to baseline models such as BLOOM and BLOOMZ. Categories of these baseline models include pre-trained large models with decoder-only architecture like BLOOM, and instruction-following language models like BLOOMZ and ChatGPT. The evaluation metric of the EcomInstruct involves converting all tasks to generative paradigms and using text generation evaluation metrics like ROUGE-L. Classification tasks were evaluated with precision, recall, and F1 scores. The EcomInstruct dataset, comprising 12 tasks across four major categories, is divided into training and testing sections. The EcomGPT is trained on 85,746 instances of E-commerce data. The performance of the model is assessed based on its capacity to generalize unseen tasks or datasets, with emphasis on cross-language and cross-task paradigm settings.


{\paragraph{Medicine}}
\citet{wang2023cmb} propose a localized medical benchmark called CMB, or the Comprehensive Medical Benchmark in Chinese. CMB is rooted entirely in the native Chinese linguistic and cultural framework. While traditional Chinese medicine is a significant part of this evaluation, it does not make up the entire benchmark.
Both prominent LLMs, such as ChatGPT and GPT-4, and localized Chinese LLMs, including those specializing in the health domain, are evaluated using CMB. However, the benchmark is not designed as a leaderboard competition, but rather as a tool for self-assessment and understanding the progression of models in this field.
CMB embodies a comprehensive, multi-layered medical benchmark in Chinese, comprised of hundreds of thousands of multiple-choice questions and complex case consultation questions. This wide range of queries covers all clinical medical specialties and various professional levels, seeking to evaluate a model's medical knowledge and clinical consultation capabilities comprehensively.
 
\citet{li2023huatuo} introduce Huatuo-26M dataset, the largest Chinese medical Question and Answer (QA) dataset to date, including over 26 million high-quality medical QA pairs. It covers a wide range of topics, such as diseases, symptoms, treatments, and drug information. The dataset is a valuable resource for anyone looking to improve AI applications in the medical field, such as chatbots and intelligent diagnostic systems.
The Huatuo-26M dataset is gathered and integrated from various sources, including online medical encyclopedias, online medical knowledge bases, and online medical consultation records. Each QA pair in the dataset contains a problem description and a corresponding answer from a doctor or expert. Although a significant proportion of the Huatuo-26M dataset is constituted by the online medical consultation records, the text format data from these records is not publicly available, for unspecified reasons.
This dataset is expected to be instrumental for multiple types of research and AI applications in the medical field. These areas of application extend to Natural Language Processing tasks like developing QA systems, text classification, sentiment analysis, and Machine Learning model training tasks like disease prediction and personalized treatment recommendation. Consequently, the dataset is favorable for developing AI applications in the medical field, from intelligent diagnosis systems to medical consultation chatbots.

% {\paragraph{Gene: }}
\citet{jin2023genegpt} use nine tasks related to NCBI resources to evaluate the proposed GeneGPT model. The dataset used is GeneTuring benchmark~\cite{Hou2023GeneTuring}, which contains 12 tasks with 50 question-answer pairs each. The tasks are divided into four modules - Nomenclature  (Gene alias, Gene name conversion),  Genomic location(Gene SNP association, Gene location, SNP location), 
  Functional analysis(Gene disease association, Protein-coding genes), 
  Sequence alignment (DNA to human genome, DNA to multiple species). Two settings of GeneGPT are assessed: a full setting where all prompt components are used, and a slim setting using only two components. The performance of GeneGPT is compared against various baselines including general-domain GPT-based LLMs like GPT-2, GPT-3, and ChatGPT. Additionally, GPT-2-sized biomedical domain-specific LLMs such as BioGPT and BioMedLM are evaluated. The new Bing, a retrieval-augmented LLM with access to relevant web pages, is also assessed. The result evaluation of the compared methods is based on the results reported in the original benchmark and are manually evaluated. However, the evaluation of the proposed GeneGPT method is determined through automatic evaluations.





\paragraph{Law}
LegalBench~\citep{guha2023legalbench} is a benchmark for legal reasoning introduced due to the increasing use of LLMs in the legal field. It consists of 162 tasks covering six types of legal reasoning and was collaboratively constructed through an interdisciplinary process with significant contributions from legal professionals. The tasks designed aim to either demonstrate practical legal reasoning capabilities or measure reasoning skills that are of interest to lawyers. To facilitate discussions between different fields relating to LLMs in law, LegalBench tasks correspond to popular legal frameworks for describing legal reasoning, thus creating a shared language between lawyers and LLM developers. The paper not only describes LegalBench, but also presents an evaluation of 20 different open-source and commercial LLMs and highlights the types of research opportunities that LegalBench can provide.

LawBench~\citep{fei2023lawbench} is an evaluation framework dedicated to assessing the capabilities of LLMs in relation to legal tasks. The context-specificity and high-stakes nature of the law field make it crucial to have a clear grasp of LLMs' legal knowledge and their ability to execute legal tasks. 
LawBench dives deep into three cognitive evaluations of LLMs: the capability to memorize crucial legal details, understanding legal texts, and applying legal knowledge to resolve complicated legal problems. A total of 20 diverse tasks have been put together, covering five main task categories: single-label classification, multi-label classification, regression, extraction, and generation. 
Throughout the evaluation process, 51 LLMs were extensively tested under LawBench, compiling a spectrum of language models including 20 multilingual LLMs, 22 Chinese-oriented LLMs, and 9 law-specific LLMs. The results concluded that GPT-4 ranked as the superior LLM in the law domain, significantly surpassing its competitors. Despite the noted improvements when LLMs were fine-tuned on law-specific texts, the study acknowledged that there is still a long road ahead in achieving highly reliable LLMs for legal tasks.

\citet{bi2023oceangpt}  design OceanBench to evaluate the capabilities of LLMs for oceanography tasks.
OceanBench includes a total of 15 type of ocean-related tasks, such as question-answering and description tasks.
The samples in OceanBench are are automatically generated from the seed dataset and have undergone manual verification by experts.
